\begin{table}[H]
\centering
\caption{Nama Daerah Perbandingan Cluster Biplot vs STATIS (Tahunan) (2022)} \label{tab:app1biplot_statis_comp_name_2022}
\begin{adjustbox}{max width=\textwidth}
\begin{tabular}{|p{4cm}|p{10cm}|}
\hline
\textbf{Kategori} & \textbf{Daftar Wilayah} \\
\hline
Biplot 1 (Rendah) \& Statis 1 (Rendah) & Kep. Seribu, Pandeglang, Lebak, Serang, Kota Cilegon \\
\hline
Biplot 1 (Rendah) \& Statis 2 (Sedang) & Tangerang, Kota Serang \\
\hline
Biplot 1 (Rendah) \& Statis 3 (Tinggi) & - \\
\hline
Biplot 2 (Sedang) \& Statis 1 (Rendah) & Sukabumi, Cianjur, Garut, Tasikmalaya, Ciamis, Majalengka, Indramayu, Pangandaran \\
\hline
Biplot 2 (Sedang) \& Statis 2 (Sedang) & Bogor, Bandung, Kuningan, Cirebon, Sumedang, Subang, Purwakarta, Karawang, Bandung Barat, Kota Cirebon, Kota Tasikmalaya, Kota Banjar \\
\hline
Biplot 2 (Sedang) \& Statis 3 (Tinggi) & - \\
\hline
Biplot 3 (Tinggi) \& Statis 1 (Rendah) & - \\
\hline
Biplot 3 (Tinggi) \& Statis 2 (Sedang) & Kota Sukabumi \\
\hline
Biplot 3 (Tinggi) \& Statis 3 (Tinggi) & Kota Jakarta Selatan, Kota Jakarta Timur, Kota Jakarta Pusat, Kota Jakarta Barat, Kota Jakarta Utara, Bekasi, Kota Bogor, Kota Bandung, Kota Bekasi, Kota Depok, Kota Cimahi, Kota Tangerang, Kota Tangerang Selatan \\
\hline
\end{tabular}
\end{adjustbox}
\end{table}